\documentclass[11pt]{article}
\usepackage[a4paper,margin=2.3cm]{geometry}
\usepackage{amsmath,amssymb,amsthm,bm}
\usepackage[dvipsnames]{xcolor}\usepackage{graphicx,booktabs,array,tabularx,caption}
\usepackage[numbers,sort&compress]{natbib}
\usepackage[colorlinks=true,linkcolor=blue!50!black,citecolor=blue!50!black,urlcolor=blue!50!black]{hyperref}
\captionsetup{font=small,labelfont=bf}
\newtheorem{theorem}{Theorem}\newtheorem{proposition}{Proposition}\newtheorem{corollary}{Corollary}
\theoremstyle{definition}\newtheorem{definition}{Definition}\theoremstyle{remark}\newtheorem{remark}{Remark}
\newcommand{\im}{\operatorname{Im}}\newcommand{\re}{\operatorname{Re}}
\linespread{1.08}

\title{Mass-metric Lagrangian Hodge--Dirac calculus on tracer trajectory complexes:\\ exact finite-step budgets for compressible transport and mixing}
\author{First Author$^{1}$ \and Second Author$^{1}$\\[0.3em]{\small $^1$Affiliation, City, Country}}
\date{Draft v2.1 (revised and verified), October 2026}

\begin{document}
\maketitle

\begin{abstract}
Lagrangian particle data are the natural record of transport in compressible and shock-dominated flows, yet most quantitative Lagrangian diagnostics require velocity gradients that shocks render ill-defined, and most network and transfer-operator methods measure the transport of area rather than of mass. We develop a calculus that uses tracer positions alone, plus the density at one reference time, and is exact at the level of finite time steps. On the Delaunay complex of the tracers, every triangle undergoes an affine three-point map $F\zeta=p\zeta+q\bar\zeta$; a complex displacement 1-cochain satisfies $(d_1c)_t=2A_t\im p_t+i\,\Delta A_t$ for any step size, so that circulation and material area change are two components of one algebraic identity, and a position-only discrete circulation identity, from which Kelvin's theorem follows, holds exactly. The identities are metric-free and supply a material mass ledger that telescopes across shocks; equipping the complex with the resulting tracer masses as metric yields the Laplacian blocks of a Dirac operator: a variable-density Lagrangian diffusion operator $L_0=\delta_1d_0$ with conductance $\rho\kappa$, and its time average, the weighted dynamic Laplacian. A Nanson identity splits interfacial mixing exactly into tangential stretching and compression. Validated on seven flows, including three-level-refined inviscid and viscous DNS of shock--bubble interaction and a single-mode Richtmyer--Meshkov flow, the method reconstructs density with median errors of at most $1.3\%$ (machine precision in the Sod problem), one to two orders of magnitude more accurately than DTFE or SPH estimates built from the same masses; mass-metric diffusion is 3.8--10.6 times more accurate than its area-metric counterpart, a gap set by the equation rather than by the density estimate; the position-only circulation of material bands on a shocked bubble agrees with the DNS velocity circulation along the same material boundary to $0.04$--$0.11\%$ (median) and with the Samtaney--Zabusky model to $8\%$; and the Richtmyer--Meshkov growth rate obtained from tracer positions agrees with the impulsive model to $3\%$, its post-shock parameters with the exact shock-refraction solution to $0.7\%$. Every algebraic identity is verified symbolically, and the ledger is verified on manufactured flow maps: it is exact across planar shocks, its error across curved shocks depends on $h/R_s$ alone, and it is first order in smooth flow, where a point-wise tensor variant is second order. A noise propagation law, an adaptive anchoring rule and a per-track ledger for broken trajectories make the method applicable to particle-tracking data.
\end{abstract}

\section{Introduction}
Shock-accelerated interfaces \citep{HaasSturtevant1987,PiconeBoris1988,SamtaneyZabusky1994}, supersonic mixing layers and shock--boundary-layer interactions transport and mix material across shocks, contact surfaces and slip lines. Lagrangian coherent structure (LCS) theory characterises transport from the deformation of the flow map \citep{Haller2015,Hadjighasem2017}, but its standard diagnostics differentiate the velocity field or the flow map, which is fragile where shocks deprive the flow of regularity. Set-oriented and spectral approaches based on transfer operators and the dynamic Laplacian \citep{Froyland2015} identify coherent sets from trajectory data without derivatives; Froyland and Junge \citep{FroylandJunge2018} discretised the dynamic Laplacian by finite elements on scattered, sparse and incomplete trajectories, Froyland and Kwok \citep{FroylandKwok2020} extended it to non-volume-preserving dynamics on weighted manifolds, and Karrasch and Keller \citep{KarraschKeller2020} showed that it is the time average of the Lagrangian heat equation. Trajectory-network methods \citep{Iacobello2021,PadbergGehle2017,BanischKoltai2017} and single-trajectory or relative diagnostics \citep{Haller2021,Aksamit2024} further reduce the data requirements. In compressible flows two difficulties remain: (i) the flow map preserves mass, not area, so that area-weighted transport, diffusion and composition are biased wherever the flow compresses; and (ii) the mass measure itself is not available from tracer positions unless density can be tracked across shocks.

This paper addresses both difficulties with a calculus on the trajectory complex that is exact at finite time steps. Its ingredients are elementary --- the affine map of three points, the Delaunay star of a tracer, the P1 finite element --- but three properties make them useful for shock-dominated flows. First, the quantities are algebraic identities of the positions, not discretisations of derivatives: circulation and area change of a material triangle over one step are exact for any step size and remain meaningful across a shock (Sec.~\ref{sec:identities}). Second, a material mass ledger telescopes across shocks, so that tracer masses, and hence the mass metric, are available from positions (Sec.~\ref{sec:ledger}). Third, with this metric variable-density diffusion and the weighted dynamic Laplacian become Laplacian blocks of one Dirac operator whose coboundaries carry the metric-free identities (Sec.~\ref{sec:dirac}), and a Nanson identity attributes interfacial mixing exactly to stretching and compression (Sec.~\ref{sec:mixing}).

Relative to the dynamic-Laplacian literature \citep{FroylandJunge2018,FroylandKwok2020,KarraschKeller2020}, the finite-element stiffness on scattered trajectories and the use of weights are not new; what is new here is (a) obtaining the weights --- the tracer masses --- from positions in flows with shocks, (b) solving the transient variable-density Lagrangian advection--diffusion problem with conductance $\rho\kappa$ and validating it against DNS scalars, (c) the exact finite-step circulation and area identities and their shock-robust block form --- the area part is the two-dimensional discrete geometric conservation law of moving-mesh methods \citep{ThomasLombard1979,Farhat2001}, and what is new is its pairing with circulation in one complex cochain and its use as a mass ledger --- and (d) quantitative mixing attribution with consistent estimators. Density estimation from point sets is classical in cosmology (DTFE \citep{Schaap2000}) and in SPH \citep{Monaghan2005}; free-Lagrange methods \citep{FreeLagrange1985} evolve compressible flows on reconnecting Delaunay meshes, and our scalar solver is, in finite-element terms, a particle finite element method \citep{Idelsohn2004} for variable-density diffusion. We compare against DTFE and SPH directly.

\section{Trajectory complex and finite-step identities}\label{sec:identities}
\subsection{Three-point maps}
Tracers $x_i(t_n)$, $n=0,\dots,S-1$, are given at snapshot times. At every snapshot we form the Delaunay triangulation peeled from its outer boundary only, $K^n$, with coboundary operators $d_0,d_1$ (signed incidence matrices, $d_1d_0=0$). Between two snapshots, the three vertices of a triangle define a unique affine map, written in complex form as
\begin{equation}
F\zeta=p\zeta+q\bar\zeta,\quad p=\tfrac12[(F_{11}+F_{22})+i(F_{21}-F_{12})],\quad q=\tfrac12[(F_{11}-F_{22})+i(F_{21}+F_{12})],
\end{equation}
so that $J=\det F=|p|^2-|q|^2$, the singular values are $|p|\pm|q|$, the polar rotation angle is $\arg p$ (the discrete counterpart of the polar rotation angle of \citep{FarazmandHaller2016}) and $\mu=q/p$ is the Beltrami coefficient. Consecutive steps compose as $(p_2,q_2)\circ(p_1,q_1)=(p_2p_1+q_2\bar q_1,\,p_2q_1+q_2\bar p_1)$. Under a change of frame $x=Q(\beta(t))y+b(t)$, $p\mapsto e^{-i(\beta^{n+1}-\beta^n)}p$ and $q\mapsto e^{-i(\beta^{n+1}+\beta^n)}q$; aligning each step by the least-squares rigid motion (Kabsch) therefore makes all step readouts frame-indifferent.

\subsection{Complex displacement cochain}
For an oriented edge $e=(a\to b)$ let $u^{\rm m}_e=\frac12(\Delta x_a+\Delta x_b)$ be the mean displacement over the step and $q^{\rm m}_e=\frac12(q_e^n+q_e^{n+1})$ the mid-step edge vector; planar vectors are identified with complex numbers and $\overline{z}$ denotes complex conjugation. Triangles are oriented counter-clockwise at $t_n$ and areas are signed, so the statements also hold for triangles that invert during a step.
\begin{theorem}[finite-step circulation and area change]\label{thm:cplx}
The complex 1-cochain $c_e=\overline{u^{\rm m}_e}\,q^{\rm m}_e$ satisfies, for every triangle and every step,
\begin{equation}
(d_1c)_t=2A_t^n\,\im p_t+i\,(A_t^{n+1}-A_t^n).
\end{equation}
More generally, for every closed polygon $\Pi$ of tracers, $\sum_{e\in\partial\Pi}\im c_e=A_\Pi^{n+1}-A_\Pi^n$, and $\sum_{e\in\partial\Pi}\re c_e=\sum_t2A_t^n\im p_t$ for every triangulation of $\Pi$.
\end{theorem}
\begin{proof}
$\re(\bar u\,q)=u\cdot q$ and $\im(\bar u\,q)=u\times q$. The shoelace area $A(s)$ of the polygon with vertices $x^n+s\Delta x$ is quadratic in $s$, so $A(1)-A(0)=A'(\frac12)$ exactly, and $A'(\frac12)=\sum_eu^{\rm m}_e\times q^{\rm m}_e$ is the outward flux of the displacement through the mid-step polygon. For a triangle write $x^{n+1}=FX+c$ with $X=x^n$: the displacement is linear in $X$, $q^{\rm m}=\frac12(I+F)Q$, the midpoint rule is exact on every edge, and Green's theorem gives $\sum_eu^{\rm m}_e\cdot q^{\rm m}_e=A^n(K_{12}-K_{21})$ with $K=(F-I)^\top\frac12(I+F)$, whose antisymmetric part is $\frac12(F^\top-F)$; hence the real part is $A^n(F_{21}-F_{12})=2A^n\im p$. For a polygon, interior edges cancel in $d_1$, so the boundary sum equals the sum over any triangulation.
\end{proof}
The imaginary part is the two-dimensional discrete geometric conservation law of arbitrary Lagrangian--Eulerian schemes, which midpoint-configuration fluxes satisfy exactly \citep{ThomasLombard1979,Farhat2001}; Theorem~\ref{thm:cplx} pairs it with the circulation in one complex number. Both parts were verified in exact symbolic arithmetic for triangles, quadrilaterals and pentagons (Sec.~\ref{sec:verif}).
The real part is the circulation used for Kelvin budgets; the imaginary part, summed over the star of a tracer, is the increment of the material ledger of Sec.~\ref{sec:ledger}. Neither requires a velocity. The classical velocity cochain $\frac12(u_a+u_b)\cdot(x_b-x_a)$ is the small-step limit of $\re c_e/\Delta t$.

\begin{theorem}[position-form discrete circulation identity]\label{thm:kelvin}
Let $\Gamma_t^{n+1/2}=\re(d_1c^{n+1/2})_t$ and $g^n_e=\frac12(\Delta^2x_a+\Delta^2x_b)\cdot\big(q_e^n+\gamma(q_e^{n+1}-2q_e^n+q_e^{n-1})\big)$ with $\Delta^2x=x^{n+1}-2x^n+x^{n-1}$ and any $\gamma\in\mathbb R$. Then, exactly for any three snapshots,
\begin{equation}
\Gamma_t^{n+1/2}-\Gamma_t^{n-1/2}=(d_1g^n)_t.
\end{equation}
\end{theorem}
\begin{proof}
With $u^\pm$ the edge-mean displacements of the two half steps and $\delta^\pm$ the corresponding changes of the edge vector, the edge contribution is $(u^+-u^-)\cdot q^n+\frac12(u^+\cdot\delta^++u^-\cdot\delta^-)$. Since $\delta^\pm=\Delta x_b^\pm-\Delta x_a^\pm$, each $u^\pm\cdot\delta^\pm=\frac12(|\Delta x_b^\pm|^2-|\Delta x_a^\pm|^2)$ is an exact difference and sums to zero around a triangle; the $\gamma$-term is likewise $\frac\gamma2(\Delta^2x_a+\Delta^2x_b)\cdot(\Delta^2x_b-\Delta^2x_a)$.
\end{proof}
Theorem~\ref{thm:kelvin} is a summation by parts in time: it holds for any three point configurations and contains no dynamics. Kelvin's theorem follows once the dynamics enter through the acceleration.
\begin{corollary}[discrete Kelvin theorem]\label{cor:kelvin}
Let the tracers follow a smooth flow map with Lagrangian acceleration $a$, so that $\Delta^2x_i=\Delta t^2a(x_i^n,t_n)+\frac{\Delta t^4}{12}D_t^2a+O(\Delta t^6)$, and let $R$ be a material polygon with perimeter $L$ and largest edge $h$. If $a=-\nabla\Psi$ (barotropic, inviscid), then for $\gamma=0$
\begin{equation}
\Big|\sum_{t\in R}\big(\Gamma_t^{n+1/2}-\Gamma_t^{n-1/2}\big)\Big|\le\frac{\Delta t^2L}{12}\Big(h^2\|\nabla^2a\|_\infty+\Delta t^2\|D_t^2a\|_\infty\Big)+O(\Delta t^6L).
\end{equation}
\end{corollary}
\begin{proof}
By Theorem~\ref{thm:kelvin} the left side is $\sum_{e\in\partial R}g_e$. Its leading term $\Delta t^2\sum_e\frac12(a_a+a_b)\cdot q_e$ is the trapezoidal rule for $\oint_{\partial R}a\cdot dx=0$; on an edge of length $\ell$ the trapezoidal error is at most $\ell^3\|\nabla^2a\|_\infty/12$, and $\sum_e\ell_e^3\le h^2L$. The $\Delta t^4$ term is bounded by $\frac{\Delta t^4}{12}L\|D_t^2a\|_\infty$.
\end{proof}
In the continuum limit $(d_1g)_t/\Delta t^2$ is the baroclinic plus viscous torque on the triangle; across a shock, $\Delta^2x$ is a finite velocity jump times $\Delta t$, so the impulsive deposition is well defined. The sum over the triangles of a material region depends only on its boundary cochain, so budgets are independent of how the interior is re-triangulated.

\subsection{Shock detection and merged blocks}\label{sec:shock}
The single-step log area ratio of a tracer star equals $\int\theta\,dt$; it is $O(\Delta t)$ in smooth flow and $-\ln\eta=O(1)$ when a shock of density ratio $\eta$ sweeps the star. Steps with $|\ln(A^\ast_{n+1}/A^\ast_n)|>\tau=\max(0.2,4\hat\sigma)$ are marked hot; for $\gamma=1.4$ this detects shocks with $M_s\gtrsim1.13$ in one step. Composing the three-point maps of consecutive hot steps into a block $(p_B,q_B)$ removes the pseudo-rotation of triangles straddling the shock: a complete planar-shock crossing is uniaxial, $F=I-(1-\eta^{-1})nn^\top$, so that $\arg p_B=0$ and
\begin{equation}
|\mu_B|=\tanh\big(\tfrac12|\ln J_B|\big),
\end{equation}
independently of the shock orientation (Fig.~\ref{fig:signature}). The signature is a characterisation, not only a necessary condition:
\begin{proposition}[uniaxial signature]\label{prop:sig}
For $F$ with positive trace, $\arg p=0$ if and only if $F$ is symmetric. For symmetric positive-definite $F$ with eigenvalues $\lambda_{1,2}$, $|\mu|=\tanh(\frac12|\ln J|)$ if and only if $(\lambda_1-1)(\lambda_2-1)=0$.
\end{proposition}
\begin{proof}
$\im p=\frac12(F_{21}-F_{12})$. With $|\mu|=|\lambda_1-\lambda_2|/(\lambda_1+\lambda_2)$ and $\tanh(\frac12|\ln J|)=|1-\lambda_1\lambda_2|/(1+\lambda_1\lambda_2)$, the numerator of the difference of the squares factorises as $-4\lambda_1\lambda_2(\lambda_1-1)(\lambda_1+1)(\lambda_2-1)(\lambda_2+1)$.
\end{proof}
Because affine maps compose exactly, the merged block of a triangle equals the direct three-point map between the end frames of the block; merging removes the error of summing per-step angles. Departures from the signature measure shock curvature and non-uniform post-shock flow, including the vorticity jump across a curved shock \citep{Hayes1957}; on a converging cylindrical manufactured shock $|\arg p_B|\approx0.16\,h/R_s$ (Sec.~\ref{sec:verif}).
A shock that needs several snapshots to cross a star (snapshot Courant number $C_s=U_s\Delta t/h\lesssim1$; weak or viscous shocks) is missed by the single-step test. We therefore also mark the steps $s_0,\dots,s_0+w-1$ ($w\le6$) when the star area drops by more than $\max(\tau,3w\,\mathrm{med}_t|R_i|)$ over the window and the least-squares affine map of the star over the window has a principal stretch equal to one within $0.1$ in $\ln$ --- the rotation-invariant form of Proposition~\ref{prop:sig}, which rejects roll-up and shear.

\begin{figure}[t]\centering\includegraphics[width=\linewidth]{fig_signature.pdf}
\caption{(a) Merged-block signatures of three shock flows against the uniaxial curve. (b,c) Pseudo-rotation of triangles crossing a shock: the sum of per-step polar angles versus the polar angle of the merged block.}\label{fig:signature}\end{figure}

\section{Material mass ledger}\label{sec:ledger}
\subsection{Ledger and telescoping}
Each tracer carries a mass $m_i=\rho_i(t_c)A^\ast_i(t_c)/3$ that is conserved over the window; its density is $\rho_i(t)=\rho_i(t_c)\,e^{-\ln J_i(t)}$ with
\begin{equation}
\ln J_i(t)=\sum_k\ln\frac{A^\ast_{i,k}(t_{k+1})}{A^\ast_{i,k}(t_k)},
\end{equation}
where $A^\ast_{i,k}$ is the area of the star of $i$ in the triangulation of its $k$-th anchor time. Re-anchoring every five snapshots keeps the stars well shaped; re-anchoring is deferred while any vertex of the star is hot; the deferral cap is twice the 99th percentile of the measured hot-run length, between 20 and 50 snapshots.
\begin{proposition}[telescoping across a shock]
Let a planar shock of constant speed separate two uniform states. If the anchor star lies entirely ahead of the shock at the anchor time and entirely behind it at the measurement time, the map between the two times is affine on the star (the crossing time is an affine function of the initial position) and the area ratio is exact, irrespective of how, and in how many snapshots, the shock crossed the star.
\end{proposition}
Without deferral the median error of the ledger in the cylindrical blast wave rises from $0.58\%$ to $10.5\%$; on planar manufactured shocks the ledger is exact to $6\times10^{-14}$ with deferral and errs by up to $41\%$ without it (Sec.~\ref{sec:verif}).

\subsection{What the ledger estimates, and a non-affinity indicator}
The ledger returns $\hat\rho_i=\rho_i(t_c)A^\ast_i(t_c)/A^\ast_i(t)$, the reference density of the tracer times the area ratio of its anchor polygon; the target is $\rho(X_i,t_c)/J(X_i,t)$. Polygon area is a quadratic form of the vertices, so for any affine prediction $\hat x_k=Fx_k(t_a)+c$ of the star vertices with residuals $e_k$,
\begin{equation}
A^\ast(\hat x+e)=\det F\,A^\ast(t_a)+\sum_k\tfrac12\,e_k\times(\hat x_{k+1}-\hat x_{k-1})+A^\ast(e)
\end{equation}
exactly. The ledger is therefore exact when the star moves affinely, and $\delta=\ln[A^\ast(t)/A^\ast(t_a)]-\ln\det F_{\rm LS}$ is the part of an increment not explained by an affine motion. Accumulated over anchor segments it gives a per-tracer indicator. It is not an error bound, but it identifies the largest $5\%$ of the errors with an area under the ROC curve of $0.97$ (Sod), $0.93$ (shock--bubble), $0.91$ (blast) and $0.89$ (cylindrical manufactured shock); it carries no information in smooth flow ($0.53$). The noise law below is the same expansion with $e$ the position noise. In smooth non-uniform flow the polygon ratio estimates $J$ at the area centroid of the star rather than at the tracer, so the ledger is first order in $h$ (observed order $1.10$ on a manufactured map). A point-wise variant, $\ln\det$ of the chained least-squares deformation gradient (quadratic fit on 12 neighbours, same deferral), is second order (observed $2.18$) and also exact across planar shocks, but produces rare outliers; we use it where it agrees with the polygon ledger to $0.1$ in $\ln J$ ($86--100\%$ of the tracers) and keep the polygon ledger for the masses.

\subsection{Adaptive anchoring}
The ledger is exact when the anchor star moves affinely. We therefore re-anchor early when the least-squares affine residual of the star vertices, normalised by $\sqrt{A^\ast}$, exceeds $\varepsilon=0.01$, while keeping the five-snapshot cap and the shock deferral. This shortens anchor segments in vortex cores without affecting regions of affine motion: the $95$th percentile of the density error drops from $13.0\%$ to $7.1\%$ in the Lamb--Oseen pair and from $3.4\%$ to $2.9\%$ in the shock--vortex DNS, and is unchanged in the blast wave and Sod problems.

\subsection{Noise propagation}
\begin{proposition}[noise law]\label{prop:noise}
For independent isotropic position noise of standard deviation $\sigma$, the star area of a tracer is the area of the polygon of its neighbours and does not depend on the tracer itself; with neighbours $x_k$ in angular order,
\begin{equation}
\operatorname{std}(\delta\ln A^\ast)=\frac{\sigma}{2A^\ast}\Big(\sum_k|x_{k+1}-x_{k-1}|^2\Big)^{1/2}=c_\ast\frac{\sigma}{h},
\end{equation}
with $c_\ast=0.816$ for a regular hexagonal star. Intermediate re-anchoring noise largely cancels because old and new stars share vertices measured in the same snapshot, so the median ledger error is $\approx0.674\sqrt2\,c_\ast\,\sigma/h$.
\end{proposition}
On a jittered lattice $c_\ast=0.777$, predicting a median of $0.74\,\sigma/h$; measured values are $0.73$--$0.79\,\sigma/h$ for 11 snapshots and $0.80$--$0.89\,\sigma/h$ for 41, and $0.75\,\sigma/h$ (Sod) and $1.0\,\sigma/h$ (shock--vortex) in flows. Unlike gradient-based diagnostics, whose sensitivity grows with stretching \citep{Badza2023,Balasuriya2020}, the ledger error is not amplified by accumulated deformation; a $1\%$ median error requires $\sigma\lesssim0.013h$, the order of typical particle-tracking accuracy.

\subsection{Broken trajectories}
Particle-tracking data contain broken tracks: a particle is lost for a few frames and reappears with a new identity. We store the data as a track table (one column per track, missing entries undefined) and let every track keep its own anchor: the ordered ring of its anchor-star neighbours and their anchor positions. When ring vertices disappear, the area ratio is evaluated on the polygon of the surviving vertices with the same vertex subset in numerator and denominator, which is exact for affine motion. When fewer than three vertices survive inside a shock stencil, the ring is rebuilt \emph{at the anchor time} from tracks present both then and now, so that the anchor stays ahead of the shock and the crossing still telescopes. New tracks start their own ledger and are calibrated from established neighbours at the first frame at which they are outside a shock stencil. Section~\ref{sec:robust} compares this per-track ledger with a ledger restricted to windows of complete tracks.

\section{Mass-metric Dirac operator}\label{sec:dirac}
Let $M_0=\operatorname{diag}(m_i)$ and let $M_1$ be the Whitney-form (P1 Galerkin) mass matrix with coefficient $(\rho\kappa)_T$ on each triangle, with $\delta_1=M_0^{-1}d_0^\top M_1$; the diagonal cotangent Hodge star gives the same $d_0^\top M_1d_0$ but a different $L_1$. The Dirac operator $D$ on $C^0\oplus C^1\oplus C^2$, with blocks $d_0,d_1$ and their adjoints, satisfies $D^2=\operatorname{diag}(L_0,L_1,L_2)$.
\begin{proposition}[transport blocks]\label{prop:blocks}
(i) The ledger increments ($\im d_1c$ summed over a star) and the Kelvin budget ($\re d_1c$) use only the metric-free coboundary $d_1$ (Theorem~\ref{thm:cplx}); the metric enters through the masses they produce. (ii) The variable-density Lagrangian diffusion $\rho\,Dc/Dt=\nabla\cdot(\rho\kappa\nabla c)$ discretised with lumped P1 elements in material coordinates is $\dot c=-L_0c$, i.e. $M\dot c=-K(t)c$ with $M=\operatorname{diag}(m_i)$ constant in time and $K_{ij}=\sum_{T\ni i,j}(\rho\kappa)_TA_T\nabla\lambda_i\cdot\nabla\lambda_j$. (iii) Its time average $\bar K=S^{-1}\sum_sK(t_s)$ with $M$ is the weighted dynamic Laplacian of \citep{FroylandKwok2020} in the finite-element form of \citep{FroylandJunge2018}; coherent sets follow from $\bar Kv=\lambda Mv$.
\end{proposition}
\begin{proposition}[conservation and dissipation]
For (ii), $\frac{d}{dt}\sum_im_ic_i=0$, and, for $\rho\kappa\ge0$ and positively oriented triangles, the mass-weighted variance decays unconditionally, \[\dot V=-2\textstyle\sum_T(\rho\kappa)_TA_T|\nabla c_T|^2\le0;\] implicit Euler steps preserve both, also when the triangulation changes between snapshots. A discrete maximum principle additionally requires $K_{ij}\le0$, which Delaunay triangulations do not guarantee when $\rho\kappa$ differs between the two triangles adjacent to an edge and one of the opposite angles is obtuse. With an edge conductance (harmonic mean of the adjacent triangles) and the cotangent weights of peeled-boundary edges clipped at zero, $K$ is an M-matrix for any conductance field.
\end{proposition}
The area metric, $M=\operatorname{diag}(A^\ast_i/3)$ with conductance $\kappa$, discretises $Dc/Dt=\kappa\nabla^2c$ and misses $\kappa\nabla\ln\rho\cdot\nabla c$; the two metrics coincide in incompressible flow. The choice of metric is therefore a physical, testable choice of measure. In the Sod and blast problems the element form never left the initial bounds by more than $3\times10^{-8}$, whereas the clipped edge form lost accuracy in the narrow Sod strip ($1.77\%$ against $0.17\%$); we keep the element form and report bound violations.

\section{Mixing attribution}\label{sec:mixing}
Let $\phi_0$ be the initial signed distance to an interface, $N=\nabla_X\phi_0/|\nabla_X\phi_0|$, $T=R^\top N$, and $\lambda=|FT|$ the tangential stretch. The two-dimensional Nanson formula $|FT|=J|F^{-\top}N|$ gives, for a passive scalar with $c_0=g(\phi_0)$ and vanishing diffusivity,
\begin{equation}\label{eq:nanson}
\ln\frac{|\nabla c|}{|\nabla c_0|}=\ln\lambda-\ln J,
\end{equation}
an exact split of gradient growth into stretching and compression. Compression is read from the ledger. Stretching is estimated by a chained tensor: the incremental deformation gradient is fitted by quadratic least squares on 12 neighbours and composed with the ledger's anchoring rule; the estimator is consistent, its error is amplified by the anisotropy $A=\sigma_{\max}(F)/\lambda$, and it converges at second order in the Lamb--Oseen pair. A covector estimate $\hat\lambda=\hat J|\widehat{\nabla_x\phi_0}|/|\nabla_X\phi_0|$ is integration-free and serves as a cross-check; both fail when the interface is a strong shear layer whose material layer the tracers do not resolve, an information limit of position data: across a manufactured shear layer of thickness $\delta$ the chained-tensor stretch error is $11\%$, $5\%$, $2\%$ and $0.7\%$ for $\delta/h=1,2,4,8$. For $\kappa>0$ the iso-lines of $c$ are not material and the gradient budget acquires a diffusive term; (\ref{eq:nanson}) is its kinematic part. It attributes the local gradient; the scalar dissipation of a thin layer scales as $\rho\kappa\lambda^2/J$, so for the integrated mixing rate stretching counts twice.

\section{Test cases}
Table~\ref{tab:cases} lists the flows. The DNS solve the two-dimensional Euler or Navier--Stokes equations with a MUSCL--HLLC scheme \citep{Toro1994}; tracers are advanced inside the solver. The shock--bubble interaction was computed on three grids ($\Delta x=0.01,0.005,0.0025$) and, with viscosity, heat conduction and species diffusion ($\mu=5\times10^{-4}$, $Re\approx1300$), on three further grids ($0.01,0.005,0.0035$), always with the same tracers.

\begin{table}[t]\centering\small\caption{Test cases and references.}\label{tab:cases}
\begin{tabularx}{\linewidth}{@{}l l X@{}}\toprule
Flow & Setting & Reference\\\midrule
Double gyre (DG) & $160\times80$ tracers, $T=10$ & analytic vorticity, exact flow map\\
Lamb--Oseen pair (LO) & $N=2000$--$32000$, $T=6\pi$ & analytic vorticity, exact flow map\\
Sod shock tube & $h=0.004$, 7500 tracers & exact solution and trajectories\\
Cylindrical blast & $h=0.022$, 9345 tracers & exact self-similar kinematics\\
Shock--isentropic vortex & $M_s=1.3$, 14400 tracers, DNS $\Delta x=0.01$--$0.0025$ & DNS at tracers\\
Shock--bubble & $M_s=1.5$, density ratio 3, 15960 tracers & inviscid and viscous DNS, three grids each\\
Richtmyer--Meshkov (RM) & $M_s=1.5$, $A=0.5$, $a_0=0.04$, $\lambda=1$, 8000 tracers & DNS, Richtmyer model\\\bottomrule
\end{tabularx}\end{table}

\section{Results}\label{sec:results}
\subsection{Identities and local readouts}
Theorems~\ref{thm:cplx} and \ref{thm:kelvin} are verified symbolically (Sec.~\ref{sec:verif}) and hold to $10^{-16}$--$10^{-13}$ in floating point on random large affine steps, non-affine perturbations, and all flow data, independently of the step size. The step vorticity $\hat\omega_t=2A^n_t\im p_t/(\bar A_t\Delta t)$ correlates with the reference vorticity at $0.97$--$0.999$ (DG), $0.93$--$0.9995$ (LO) and $0.976$--$0.9999$ (shock--vortex, excluding triangles that straddle the shock); including them lowers the correlation to $0.47$ at the time the shock crosses the vortex, which is why only integrated quantities are reported across shocks.

\begin{figure}[p]\centering\includegraphics[width=0.92\linewidth]{viz_en_dg.pdf}
\caption{Double gyre, all fields from tracer positions. Rows: material labels and the tracked material line; network diffusion scalar (mass metric); step vorticity $\hat\omega$; velocity from position differences, stretch field $\ln\lambda$ of the chained tensor, anisotropy $A$, and the ledger compression field (zero for incompressible flow).}\label{fig:vizdg}\end{figure}
\subsection{Density}
Figures~\ref{fig:vizdg}, \ref{fig:vizsb} and \ref{fig:vizrm} show the fields the method produces from positions in three of the flows; all panels except those labelled DNS use tracer positions and the initial density only.

\emph{Double gyre (Fig.~\ref{fig:vizdg}).} The material-label rows display lobe transport directly: the fraction of left-gyre material found in the right half grows as $0$, $14.6\%$, $20.0\%$ and $32.9\%$ at $t=0,3.4,6.6,10$, one lobe per forcing period, and the tracked material line $x_0=0.5$ is folded into an S-shaped filament. The network diffusion scalar is stretched along the same filaments; the share of tracers in the mixed range $0.1<c<0.9$ rises from $11\%$ to $59\%$ and the mass-weighted variance drops to $0.56$ of its initial value, with most of the decay in the second half of the window as the filaments thin (Batchelor mixing). The stretch field at $t=10$ has median $0.96$ and a $95$th percentile of $3.1$, with $23\%$ of the tracers beyond $|\ln\lambda|=2$; the anisotropy is anticorrelated with $|\ln\lambda|$ (rank correlation $-0.45$): iso-lines are stretched most where they are aligned with the maximal stretching direction ($A\approx1$). The compression field, which must vanish in this flow, stays below $0.06$ (95th percentile).

Table~\ref{tab:density} compares the ledger with DTFE and an adaptive SPH summation using the same tracer masses (Fig.~\ref{fig:baselines}a). The ledger's median error is at most $1.3\%$ in all seven flows and at machine precision in Sod --- which only says that most Sod tracers sit where the star moves affinely (uniform states and the rarefaction fan), hence the 95th percentiles --- one order of magnitude below SPH and one to two orders below DTFE, whose masses are assigned on the initial stars while the current stars have changed neighbours. The $95$th percentiles of all methods are dominated by tracers whose stars are being swept by a shock or straddle an interface at the evaluation time. With multi-window shock detection and the hybrid point-wise estimator of Sec.~\ref{sec:ledger} the median error falls to $0.11\%$ (DG), $0.29\%$ (blast) and $0.40\%$ (shock--bubble), and the Sod 95th percentile to $0.29\%$; in the rerun shock--vortex, Richtmyer--Meshkov and viscous shock--bubble DNS the median falls from $0.34\%$, $0.73\%$ and $0.53\%$ to $0.20\%$, $0.47\%$ and $0.31\%$, and the shock--vortex 95th percentile from $3.4\%$ to $1.9\%$ (the non-affine vortex-core error is what the point-wise estimator removes); the 95th percentiles at shocks are unchanged.

\begin{table}[t]\centering\small\caption{Density from positions: median / 95th percentile of the relative error at the final time.}\label{tab:density}
\begin{tabular}{lccc}\toprule
Flow & ledger & SPH & DTFE\\\midrule
DG & 1.3\% / 5.7\% & 11.4\% / 65\% & 32\% / 254\%\\
LO & 0.41\% / 13\% & 6.4\% / 49\% & 26\% / 168\%\\
Sod & $<10^{-13}$ / 1.0\% & 1.9\% / 7.1\% & 0.20\% / 38\%\\
Blast & 0.58\% / 32\% & 6.8\% / 39\% & 20\% / 107\%\\
Shock--vortex & 0.34\% / 3.4\% & 3.4\% / 32\% & 17\% / 98\%\\
Shock--bubble & 0.67\% / 30\% & 7.1\% / 44\% & 26\% / 142\%\\
RM & 0.73\% / 35\% & 4.6\% / 41\% & 17\% / 73\%\\\bottomrule
\end{tabular}\end{table}

\begin{figure}[t]\centering\includegraphics[width=\linewidth]{fig_baselines.pdf}
\caption{(a) Median density error of the ledger, SPH and DTFE with identical tracer masses. (b) Final-time diffusion error of the mass-metric network operator, SPH conduction and the area-metric network operator.}\label{fig:baselines}\end{figure}

\subsection{Diffusion}\label{sec:diff}
Errors are RMS deviations from the reference divided by the standard deviation of the reference. With the same ledger densities and implicit time step, the mass-metric network operator gives $0.17\%$ (Sod), $2.2\%$ (blast), $3.2\%$ (DG) and $6.3\%$ (shock--bubble, against the DNS scalar). The area metric is 10.6, 10.2 and 3.8 times less accurate in the compressible flows ($1.8\%$, $22\%$, $24\%$) and equivalent in the incompressible DG, as expected. The gap is set by the equation, not by the density estimate: with the exact density in the conductance the error is $0.16\%$ (Sod) and $1.85\%$ (blast), with SPH densities $0.19\%$ and $2.48\%$, with DTFE densities $0.25\%$ and $9.12\%$. Four sub-steps per snapshot interval or Crank--Nicolson change the Sod error by at most $0.03$ percentage points, so the snapshot interval is not the limiting error. The gap persists for a gas-kinetic conductance $\rho D\propto T^{0.7}$ with $T$ reconstructed from positions --- the entropy is constant along a tracer except at detected shocks, where the measured density ratio gives the Rankine--Hugoniot entropy jump ($T$ exact at the median, 95th percentile $2.8\%$): $0.126\%$ for the mass metric ($0.126\%$ with the exact $T$) against $2.9\%$ for the area metric. SPH conduction in conservative Brookshaw--Cleary--Monaghan form \citep{ClearyMonaghan1999} gives $0.88\%$, $4.8\%$, $17\%$ and $6.7\%$; in the double gyre, folded filaments make the tracer distribution strongly anisotropic, which isotropic kernels resolve poorly while the P1 stiffness remains consistent.

\subsection{Circulation and comparison with deposition models}\label{sec:circ}
Position-form region circulations agree with the DNS integrated over the same triangles to $0.16\%$ (RM), $0.8\%$ (shock--vortex), $0.5\%$ (LO) and $2.5\%$ (DG). For the shock--bubble interaction (Fig.~\ref{fig:sz}), the circulation of the material band $|\phi_0|<0.15$ of the upper half bubble is $0.266$ at $t=0.6$. The like-for-like reference, the DNS velocity integrated along the same material boundary, agrees to $0.04\%$, $0.07\%$ and $0.11\%$ for the bands $|\phi_0|<0.15$, $0.08$ and $0.04$ (median over $t>0.4$; at most $1.1\%$); the DNS vorticity summed over the whole upper half plane, $0.270$, refers to a different region. The SZ1 model of Samtaney and Zabusky \citep{SamtaneyZabusky1994} gives $0.289$, with $\Gamma=(1+\pi/2)\Gamma'_1r_0$ with $\Gamma_1'=\frac{2}{\gamma+1}(1-\eta^{-1/2})(1+M^{-1}+2M^{-2})(M-1)c_0$. In the inviscid flow the circulation of the wide band is then nearly conserved ($0.266\to0.274$), whereas a vertex-averaged quadrature of the DNS vorticity over the same, by then strongly deformed, triangles drifts to $0.43$ while the line integral stays at $0.274$. The narrow band $|\phi_0|<0.04$ loses circulation ($0.257\to0.159$) in position form and line integral alike: its boundary lies where the inviscid DNS dissipates the sheet numerically, so it is the assumption of Corollary~\ref{cor:kelvin}, not the measurement, that fails there. In the viscous flow both agree throughout (median difference $0.35\%$), and a narrow band loses circulation by viscous diffusion while a wide band does not.

\begin{figure}[t]\centering\includegraphics[width=0.55\linewidth]{fig_sz.pdf}
\caption{Circulation of the upper half of the shocked bubble: position form (material band), DNS domain integral and the Samtaney--Zabusky model.}\label{fig:sz}\end{figure}

\subsection{Mixing attribution}
In the exact flows the chained tensor reproduces the interface stretch to within $0.007$ (DG $1.708$ vs.\ $1.711$, blast $0.576$ vs.\ $0.577$, Sod $0$), and the ledger compression to within $10^{-3}$. In the shock--bubble interaction, after the shock the interface is first shortened tangentially: the stretch reaches about $-0.4$ at $t\approx0.6$ and cancels most of the compression ($\approx0.48$); this early-time conclusion is insensitive to the DNS grid. At $t=1.4$ the inviscid result depends on resolution (Sec.~\ref{sec:dnsref}); in the viscous flow ($Re\approx1300$) the stretch at $t=1.4$ is $-0.02$ to $-0.04$ and the compression $0.31$--$0.33$ on all three grids, i.e. gradient growth is essentially all compression because viscosity suppresses the roll-up.

\begin{figure}[p]\centering\includegraphics[width=0.9\linewidth]{viz_en_sb.pdf}
\caption{Shock--bubble interaction (inviscid DNS, $\Delta x=0.005$). Rows: heavy-gas tracers and interface; ledger density; DNS density; network diffusion scalar (mass metric) and the DNS scalar with the same diffusivity; stretch field, compression field, anisotropy and step vorticity at the final time. Only the DNS rows use the solver fields.}\label{fig:vizsb}\end{figure}
\emph{Shock--bubble interaction (Fig.~\ref{fig:vizsb}).} The heavy-gas tracers show the bubble first flattened by the shock (streamwise extent $[0.52,0.98]\to[0.71,0.98]$ at $t=0.6$) and then rolled into a vortex pair and convected downstream ($[1.07,1.42]$ at $t=1.4$). The ledger density follows the DNS density frame by frame with median differences of $0.3$--$0.7\%$; $87\%$ of the largest $5\%$ differences lie within $0.03$ of the interface, where the DNS contact is numerically diffused. The network diffusion scalar and the DNS scalar have the same mixed fraction ($24\%$, $33\%$, $40\%$ at $t=0.6,1.0,1.4$) and differ by $1.4$--$7.3\%$ in $L^2$. The two mixing fields separate the mechanisms spatially: compression is nearly uniform ($0.63$ in the heavy gas, $0.58$ in the light gas) and was imposed by the shock, whereas stretching is confined to the spiral arms and has both signs (interface $5$th--$95$th percentiles $-1.5$ and $1.5$): arms being wound up are lengthened, interface segments roughly parallel to the shock are shortened. The step vorticity is predominantly negative (clockwise in the upper half), consistent with the deposited circulation of Fig.~\ref{fig:sz}; the anisotropy is largest on the vortex sheet (median $5.0$ against $2.1$ away from it), which is where the chained-tensor error is amplified.

\subsection{DNS refinement: what the reference can and cannot tell}\label{sec:dnsref}
Figure~\ref{fig:dnsref} compares ledger compression with the DNS density ratio in bands of initial distance to the bubble interface. Before roll-up ($t=0.6$) the DNS values converge towards the ledger from opposite sides of the interface while the ledger is grid-independent; the $95$th percentile of the density error falls from $13.5\%$ to $10.4\%$ to $6.1\%$, and from $5.1\%$ to $1.0\%$ away from the interface. The earlier disagreement at the interface is therefore numerical diffusion of the contact in the DNS. After roll-up ($t=1.4$) the bands away from the interface still converge, but within $|\phi_0|<0.015$ neither the inviscid DNS nor the ledger converges: Kelvin--Helmholtz structures on the vortex sheet keep refining with the grid and become thinner than the tracer spacing, which is the information limit of Sec.~\ref{sec:mixing}. With physical viscosity the interface band converges for both (ledger $0.333,0.312,0.322$; DNS $0.352,0.310,0.306$). In the shock--vortex DNS the median error falls with refinement ($0.54\%\to0.27\%$) while the $95$th percentile stays at $3.4$--$3.9\%$ and lies in the vortex core, where the star motion over an anchor segment is non-affine; this part of the error is the ledger's and is what adaptive anchoring reduces.

\begin{figure}[t]\centering\includegraphics[width=\linewidth]{fig_dnsref.pdf}
\caption{Ledger compression (solid) and DNS density ratio (dashed) in bands of initial distance to the bubble interface, for three grids (light to dark). (a,b) Inviscid DNS at $t=0.6$ and $1.4$; (c) viscous DNS at $t=1.4$.}\label{fig:dnsref}\end{figure}

\subsection{Robustness}\label{sec:robust}
\paragraph{Parameters.} The hot-step threshold ($\tau_{\min}=0.1$--$0.3$) changes the median density error by at most $0.013$ percentage points in Sod, blast and shock--vortex. The anchor period is a trade-off: a three-snapshot period raises the blast-wave median from $0.58\%$ to $0.77\%$ and its 95th percentile from $32\%$ to $46\%$, whereas in the shock--vortex core the 95th percentile is $2.9\%$, $3.4\%$ and $4.7\%$ for 3, 5 and 8 snapshots; adaptive anchoring resolves the trade-off. The shock deferral is the most sensitive parameter (Sec.~\ref{sec:ledger}). The chained tensor is insensitive to the neighbour count and anchoring ($\pm2\times10^{-4}$ in the blast-ring stretch).
\paragraph{Noise and sparsity.} With $0.02h$ noise the median error is $1.5\%$ (Sod) and $1.9\%$ (shock--vortex), consistent with Proposition~\ref{prop:noise}; the hybrid point-wise estimator lowers the median to $0.94\%$ (Sod) and $2.3\%$ (shock--bubble, polygon ledger $3.1\%$). Keeping one tracer in four or one in sixteen gives median errors of $0.85\%$ and $1.8\%$ in the shock--vortex flow.
\paragraph{Broken trajectories.} Figure~\ref{fig:breaks} compares the per-track ledger with a ledger restricted to five-frame windows of complete tracks. The per-track ledger provides densities for all tracks present at the final time (coverage $100\%$, against $78$--$95\%$). In Sod its $95$th percentile is $1.3\%$, $1.6\%$ and $2.4\%$ for $1\%$, $3\%$ and $5\%$ breaks per frame (windowed: $3.1\%$, $4.9\%$, $6.6\%$), and $1.9\%$ (windowed $5.6\%$) when $30\%$ of the particles break exactly while crossing the shock. In the shock--vortex flow the two ledgers are equivalent on the tracks covered by both ($95$th percentile $4.2\%$), and the per-track ledger additionally covers the $14\%$ of tracks the windowed ledger loses. In the blast wave, new tracks born in the thin, dense shocked shell are calibrated across steep density gradients; the per-track median error ($1.1$--$6.5\%$) is much lower than the windowed one ($11$--$15\%$), but the $95$th percentile remains $38$--$56\%$.

\begin{figure}[t]\centering\includegraphics[width=\linewidth]{fig_breaks.pdf}
\caption{Broken trajectories: $95$th percentile of the final-time density error for the windowed and the per-track ledger; numbers above the grey bars are the windowed coverage (per-track coverage is $100\%$); dashed lines are the ledger without breaks.}\label{fig:breaks}\end{figure}

\subsection{Verification}\label{sec:verif}
\paragraph{Exact identities and invariances.} Theorems~\ref{thm:cplx} and \ref{thm:kelvin} (for polygons), the composition and frame rules of $(p,q)$, Proposition~\ref{prop:sig}, the two-dimensional Nanson identity, the exactness of reduced rings, the area expansion of Sec.~\ref{sec:ledger} and $c_\ast=\sqrt{2/3}$ were verified in exact symbolic arithmetic. Rigid motions with time-dependent rotation, permutations of the tracers and uniform scaling change the Sod ledger by at most $1.5\times10^{-12}$; running the ledger backwards from the final frame recovers the initial density at machine precision (median) and to $1.1\%$ (95th percentile).
\paragraph{Manufactured flow maps.} To separate the method from the DNS we prescribe analytic flow maps with exact density (Fig.~\ref{fig:verif}). (a) Planar shocks with $\theta\in\{0,0.5\}$, $\eta\in\{1.5,3,6\}$, tangential slip and $C_s=0.5$--$4$: whenever the shock is detected the ledger is exact to $6\times10^{-14}$ and the merged block is uniaxial to $7\times10^{-13}$; without deferral the error reaches $41\%$. The single-step test misses weak or slow shocks ($\eta=1.5$, $C_s=0.5$: maximum error $12\%$; $\eta=3$, $C_s=0.25$: $58\%$); the multi-window test of Sec.~\ref{sec:shock} restores exactness in both. (b) A converging cylindrical shock with a sink flow behind it: the error depends on $h/R_s$ alone (curves for $h=0.04,0.02,0.01$ collapse), with medians of $0.08--0.34\%$ for $h/R_s=0.01$--$0.12$. (c) A smooth compressible, rotational map: observed order $1.10$ for the polygon ledger and $2.18$ for the tensor ledger. (f) A thin shear layer quantifies the information limit of the stretch estimate.
\paragraph{Like-for-like references.} (d) Circulation against the DNS line integral along the same boundary (Sec.~\ref{sec:circ}); in the rerun shock--vortex core ($r_0<2R_v$), Richtmyer--Meshkov interface band and viscous shock--bubble bands the median differences are $0.04\%$, $0.03\%$ and $0.03$--$0.04\%$, whereas the vertex-vorticity quadrature differs by up to $1.8\%$ (median); (e) the diffusion decomposition of Sec.~\ref{sec:diff}. The inviscid shock--bubble DNS was rerun from the archived script. Non-finite star-area ratios, which earlier code replaced silently by the previous value, are now counted; none occurred in any case.

\begin{figure}[t]\centering\includegraphics[width=\linewidth]{fig_verification.pdf}
\caption{Verification. (a) Planar manufactured shocks: maximum density error of the ledger. (b) Converging cylindrical shock: median ledger error and block rotation versus $h/R_s$ for three tracer spacings. (c) Smooth manufactured map: polygon versus point-wise tensor ledger. (d) Shock--bubble circulation of two material bands: position form, DNS line integral along the same boundary, and the vertex-vorticity quadrature used previously. (e) Diffusion error with exact, ledger, SPH and DTFE densities in the conductance, and with the area metric. (f) Chained-tensor stretch error across a manufactured shear layer of thickness $\delta$.}\label{fig:verif}\end{figure}

\paragraph{Coherent sets.} With the ledger masses, the weighted dynamic Laplacian identifies the Lamb--Oseen cores (best Jaccard index $0.82$ for the leading eigenvectors against the union of the two cores), where a mass-weighted Ulam transfer matrix on the same data gave $0.11$.

\section{Application: single-mode Richtmyer--Meshkov instability}
A $M_s=1.5$ shock passes from light gas into a heavy gas ($A=0.5$) across $x=0.5+a_0\cos(2\pi y)$ in a half-wavelength channel. Only tracer positions and the initial density are used (Fig.~\ref{fig:rm}). The shock is detected for all tracers without false alarms; density is reconstructed with median errors of $0.5$--$0.7\%$. The amplitude measured from the interface tracers on the symmetry lines agrees with the $Y=0.5$ contour of the DNS to within $1.2\%$ at five times. After the shock the amplitude is compressed to $a_0^+=0.0275$ and grows linearly at $0.0423$ for $t\in[0.6,1.4]$, before nonlinear saturation reduces it to $0.0389$. Richtmyer's impulsive model \citep{Richtmyer1960}, $\dot a=k\,\Delta u\,A^+a_0^+$, evaluated entirely from positions ($\Delta u=0.510$ from the interface translation, $A^+=0.49$ from ledger densities $2.17$ and $6.37$), gives $0.0433$, within $3\%$; the position-derived parameters agree with the exact one-dimensional shock-refraction solution ($\Delta u=0.5112$, post-shock densities $2.178$ and $6.413$, $A^+=0.493$) to within $0.7\%$. Interpreting the circulation deposited on the interface band, measured from positions to $0.16\%$ of the DNS, as a sinusoidal vortex sheet gives $\dot a=k\Gamma/4=0.050$, $19\%$ high, which reflects the thin-sheet idealisation rather than the measurement. The heavy-gas mass carried behind the mean interface agrees with the DNS to within $\pm4\%$ after the first output time. The mixing attribution shows that up to $ka\approx0.6$ the interfacial gradient growth is $\approx94\%$ compression ($-\ln J\approx0.80$, acquired at shock passage) and only slowly acquires stretch ($\ln\lambda=0.054$ at $t=2$); the growth of $\ln\langle\lambda\rangle$ matches the growth of the DNS interface length ($0.077$ vs.\ $0.076$).

\begin{figure}[p]\centering\includegraphics[width=0.92\linewidth]{viz_en_rm.pdf}
\caption{Single-mode Richtmyer--Meshkov from tracer positions: heavy-gas tracers and interface, ledger density, step vorticity (the baroclinic vortex sheet on the interface), and at $t=2$ the density error against the DNS, the stretch and compression fields, and the post-shock velocity.}\label{fig:vizrm}\end{figure}
Figure~\ref{fig:vizrm} shows the same flow as fields. The heavy-gas spike penetrates the light gas on the symmetry line $y=0$ and the light-gas bubble the heavy gas at $y=0.5$; $65\%$ of the integrated $|\hat\omega|$ is carried by tracers within $0.03$ of the interface, i.e. the vorticity is the baroclinic vortex sheet that drives the growth. At $t=2$ the stretch is positive at the spike tip ($+0.34$) and the bubble tip ($+0.07$) and negative on the flanks ($-0.12$), but its magnitude is still small compared with the shock compression ($0.80$ on the interface, $0.76$--$0.78$ in both gases). The apparent density error of $35\%$ in a band $|\phi_0|<0.02$ is the numerical diffusion of the DNS contact, as in the shock--bubble refinement; away from the interface ledger and DNS agree to $0.44\%$ (median).

\begin{figure}[t]\centering\includegraphics[width=\linewidth]{fig_rm.pdf}
\caption{Single-mode RM from tracer positions. (a) Amplitude from interface tracers and from the DNS, with the slopes of the Richtmyer model and of the circulation-based estimate; (b) circulation of the interface band; (c) mixing attribution.}\label{fig:rm}\end{figure}

\section{Discussion}
The calculus replaces derivatives by finite-step algebra on material triangles. This has three consequences that matter for compressible flows. First, quantities stay meaningful across shocks: circulation and area change of a triangle are exact for one step however large the jump, and merged blocks make the polar rotation and the shock signature well defined. Second, the metric of the complex can be chosen to be the physical measure, which the ledger supplies from positions; once it is, transport, diffusion and coherent-set computations become mass-consistent, and the difference from area-weighted versions is large and measurable. Third, errors have identifiable sources: the noise law, non-affine motion within an anchor star, and the information limit of material layers thinner than the tracer spacing. The DNS refinements show that the method's comparisons with DNS are limited by the reference as much as by the method: before roll-up the ledger is closer to the converged solution than the DNS at practical resolution.

Limitations are equally specific. All results are two-dimensional; the identities extend to tetrahedra and the three-dimensional Nanson formula, but no three-dimensional case has been computed. A reference density is required at one time. The polygon ledger is first order in smooth flow; the second-order point-wise variant needs the consistency check against it. The multi-window shock test restores exactness on manufactured weak and slow shocks but has not been tested on resolved viscous shocks. Conclusions inside rolled-up interfaces depend on the Reynolds number and on tracer resolution. New tracks born inside thin shocked shells are hard to calibrate. Real particle-tracking data, with their correlated errors and seeding inhomogeneities, have not yet been analysed.

\section{Conclusions}
Tracer positions alone determine, exactly at the level of finite steps, the circulation and area change of every material triangle; with the tracer masses obtained from a shock-telescoping ledger as metric, the same algebraic objects yield mass transport, variable-density diffusion, coherent sets, circulation budgets and an exact stretching--compression decomposition of mixing. On seven flows with analytic, exact and DNS references, density is recovered one to two orders of magnitude more accurately than by DTFE or SPH, mass-metric diffusion is 3.8--10.6 times more accurate than area-metric diffusion, also for $\rho D\propto T^{0.7}$, deposited circulation agrees with the DNS circulation along the same material boundary to $0.1\%$ and with the Samtaney--Zabusky model to $8\%$, and the Richtmyer--Meshkov growth rate follows from tracer positions to within $3\%$ of the impulsive model. The method is robust to realistic position noise, sparsity and broken tracks, and its failure modes are identified rather than hidden.

\paragraph{Data and code.} The code, the scripts that produce every number in this paper, symbolic and metamorphic tests, the manufactured flow maps, a script that checks every quoted number against the result files, and a runnable notebook are provided with the submission.

\begin{thebibliography}{99}\small
\bibitem{Haller2015} G. Haller, Lagrangian coherent structures, Annu. Rev. Fluid Mech. 47, 137--162 (2015).
\bibitem{Hadjighasem2017} A. Hadjighasem, M. Farazmand, D. Blazevski, G. Froyland, G. Haller, A critical comparison of Lagrangian methods for coherent structure detection, Chaos 27, 053104 (2017).
\bibitem{Froyland2015} G. Froyland, Dynamic isoperimetry and the geometry of Lagrangian coherent structures, Nonlinearity 28, 3587--3622 (2015).
\bibitem{FroylandJunge2018} G. Froyland, O. Junge, Robust FEM-based extraction of finite-time coherent sets using scattered, sparse, and incomplete trajectories, SIAM J. Appl. Dyn. Syst. 17, 1891--1924 (2018).
\bibitem{FroylandKwok2020} G. Froyland, E. Kwok, A dynamic Laplacian for identifying Lagrangian coherent structures on weighted Riemannian manifolds, J. Nonlinear Sci. 30, 1889--1971 (2020).
\bibitem{KarraschKeller2020} D. Karrasch, J. Keller, A geometric heat-flow theory of Lagrangian coherent structures, J. Nonlinear Sci. 30, 1849--1888 (2020).
\bibitem{Iacobello2021} G. Iacobello, L. Ridolfi, S. Scarsoglio, A review on turbulent and vortical flow analyses via complex networks, Physica A 563, 125476 (2021).
\bibitem{PadbergGehle2017} K. Padberg-Gehle, C. Schneide, Network-based study of Lagrangian transport and mixing, Nonlinear Process. Geophys. 24, 661--671 (2017).
\bibitem{BanischKoltai2017} R. Banisch, P. Koltai, Understanding the geometry of transport: diffusion maps for Lagrangian trajectory data unravel coherent sets, Chaos 27, 035804 (2017).
\bibitem{Haller2021} G. Haller, N. Aksamit, A. P. Encinas-Bartos, Quasi-objective coherent structure diagnostics from single trajectories, Chaos 31, 043131 (2021).
\bibitem{Aksamit2024} N. O. Aksamit, A. P. Encinas-Bartos, G. Haller, D. E. Rival, Relative fluid stretching and rotation for sparse trajectory observations, J. Fluid Mech. 996, A40 (2024).
\bibitem{Schaap2000} W. E. Schaap, R. van de Weygaert, Continuous fields and discrete samples: reconstruction through Delaunay tessellations, Astron. Astrophys. 363, L29--L32 (2000).
\bibitem{Monaghan2005} J. J. Monaghan, Smoothed particle hydrodynamics, Rep. Prog. Phys. 68, 1703--1759 (2005).
\bibitem{FreeLagrange1985} M. J. Fritts, W. P. Crowley, H. Trease (eds.), The Free-Lagrange Method, Lecture Notes in Physics 238, Springer (1985).
\bibitem{FarazmandHaller2016} M. Farazmand, G. Haller, Polar rotation angle identifies elliptic islands in unsteady dynamical systems, Physica D 315, 1--12 (2016).
\bibitem{Badza2023} A. Badza, T. W. Mattner, S. Balasuriya, How sensitive are Lagrangian coherent structures to uncertainties in data?, Physica D 444, 133580 (2023).
\bibitem{Balasuriya2020} S. Balasuriya, Stochastic sensitivity: a computable Lagrangian uncertainty measure for unsteady flows, SIAM Rev. 62, 781--816 (2020).
\bibitem{ClearyMonaghan1999} P. W. Cleary, J. J. Monaghan, Conduction modelling using smoothed particle hydrodynamics, J. Comput. Phys. 148, 227--264 (1999).
\bibitem{Toro1994} E. F. Toro, M. Spruce, W. Speares, Restoration of the contact surface in the HLL-Riemann solver, Shock Waves 4, 25--34 (1994).
\bibitem{SamtaneyZabusky1994} R. Samtaney, N. J. Zabusky, Circulation deposition on shock-accelerated planar and curved density-stratified interfaces: models and scaling laws, J. Fluid Mech. 269, 45--78 (1994).
\bibitem{PiconeBoris1988} J. M. Picone, J. P. Boris, Vorticity generation by shock propagation through bubbles in a gas, J. Fluid Mech. 189, 23--51 (1988).
\bibitem{HaasSturtevant1987} J.-F. Haas, B. Sturtevant, Interaction of weak shock waves with cylindrical and spherical gas inhomogeneities, J. Fluid Mech. 181, 41--76 (1987).
\bibitem{Richtmyer1960} R. D. Richtmyer, Taylor instability in shock acceleration of compressible fluids, Commun. Pure Appl. Math. 13, 297--319 (1960).
\bibitem{ThomasLombard1979} P. D. Thomas, C. K. Lombard, Geometric conservation law and its application to flow computations on moving grids, AIAA J. 17, 1030--1037 (1979).
\bibitem{Farhat2001} C. Farhat, P. Geuzaine, C. Grandmont, The discrete geometric conservation law and the nonlinear stability of ALE schemes for the solution of flow problems on moving grids, J. Comput. Phys. 174, 669--694 (2001).
\bibitem{Hayes1957} W. D. Hayes, The vorticity jump across a gasdynamic discontinuity, J. Fluid Mech. 2, 595--600 (1957).
\bibitem{Idelsohn2004} S. R. Idelsohn, E. O\~nate, F. Del Pin, The particle finite element method: a powerful tool to solve incompressible flows with free-surfaces and breaking waves, Int. J. Numer. Methods Eng. 61, 964--989 (2004).
\end{thebibliography}
\end{document}
